\section{立案管辖}
\begin{enumerate}
  \item \kd{普通刑事案件}都是\zd{公安}来行使管辖权的。
        \begin{enumerate}
          \item 犯罪地为主，嫌疑人居住地为辅。
          \item 默认县级。
          \item 市以上。\zd{涉外}犯罪、\zd{恐怖}活动犯罪、\zd{集团}犯罪、\zd{跨区}域犯罪、\zd{危害国家}安全、\zd{经济}犯罪。
        \end{enumerate}
  \item \kd{检察院自侦案件}。
        \begin{enumerate}
          \item 司法职权类。\zd{监委共管}。司法工作人员，职权，司法公正。\zd{市检 \hz 市检交下级}。
                \begin{enumerate}
                  \item \bs{搜拘滥私逼，暴虐枉玩执，放假监脱逃}。
                  \item 非法搜查、非法拘禁、滥用职权、徇私枉法、刑讯逼供。
                  \item 暴力取证、虐待被监管人、民事行政枉法裁判、玩忽职守、执行判决裁定失职/滥用职权。
                  \item 私放在押人员、徇私舞弊减刑假释、监外执行、失职致使在押人员脱逃。
                \end{enumerate}
          \item \zd{非司法职权类}。公安共管。\bs{公}安管辖的国家\bs{机}关工作人员利用职\bs{权}的重\bs{大}犯罪。\bs{省以上}检察院决定 \yjt 交办市检。
        \end{enumerate}
  \item \kd{监委}。身份标准 + 罪名标准。一定是\bs{公职人员的职务犯罪}。
        \begin{enumerate}
          \item 公务员、法律法规授权、国企、公办科教文卫体、基层自治\zd{管理}人员、其他\zd{履行公职的人员}。
          \item \bs{贪}污贿赂、\bs{玩}忽职守、\bs{滥}用职权、徇私\bs{舞}弊、重大\bs{责}任事故、\bs{其他}。
          \item 行贿受贿。
        \end{enumerate}
  \item \kd{法院}。唯一没有侦查权的机关。\bs{亲告、轻微、公转自}。
        \begin{enumerate}
          \item 亲告。\zd{侮}辱、\zd{诽}谤、\zd{暴}力干涉婚姻自由、\zd{虐}待、\zd{侵}占。\bs{只有侵占是绝对的亲告}。
          \item 轻微。\zd{重}婚、故意伤害\zd{轻}伤、\zd{伪}劣、遗\zd{弃}、\zd{通}信自由、\zd{知}识产权、非法侵入\zd{住}宅、\zd{3 下。可公可自}。
          \item 公转自。证据证明遭侵犯、曾经提出过控告、公检却不予追究。
        \end{enumerate}
  \item 国安。
  \item \kd{竞合}。原则\bs{该谁给谁}。\al{公安在查公安甲的强奸，发现还有贪污、刑讯逼供、侵占、重婚}。
        \begin{enumerate}
          \item 公自交叉。该谁给谁。\al{侵占告知被告另诉。重婚可一并查}。% https://q.uiver.app/?q=WzAsNixbMCwxLCLlvLrlpbgiXSxbMSwxLCLoh6ror4kiXSxbMiwwLCLkvrXljaAiXSxbMiwyLCLph43lqZoiXSxbMywyLCJcXGJ1bGxldCJdLFszLDAsIlxcYnVsbGV0Il0sWzAsMSwi5Y+R546w5pivIl0sWzEsMl0sWzEsM10sWzMsNCwi5bm25p+lIiwyXSxbMiw1LCLlkYrnn6UiXV0=
\[
  \begin{tikzcd}[mytest, row sep=tiny]
    &&& \txt{侵占罪} && {} \\
    \txt{侦查强奸罪} && \txt{自诉案件} \\
    &&& \txt{重婚罪} && {}
    \arrow["\txt{发现还有}", from=2-1, to=2-3]
    \arrow[from=2-3, to=1-4]
    \arrow[from=2-3, to=3-4]
    \arrow["\txt{一并查}"', no head, from=3-4, to=3-6]
    \arrow["\txt{告知另诉}", no head, from=1-4, to=1-6]
  \end{tikzcd}
\]

          \item 自公交叉。该谁给谁。这里的自指法院审理自诉案件。
          \item 公检交叉。该谁给谁，再看主罪（量刑重的）。\al{刑讯逼供给检。If(强奸)重，公安主侦查，检协助}。
          \item \bs{普通罪名和监委交叉。该谁给谁，监察为主}。该普通罪名是司法职权类 \xjt。\al{刑讯逼供、贪污受贿}。
                \begin{enumerate}
                  \item 监检竞合，同级沟通；监委合适，全案移送。
                  \item 可监可检，分别管辖；沟通上报，\zd{不停侦查}。
                \end{enumerate}
        \end{enumerate}
  \item \kd{并案管辖}。可以\zd{在职责范围内并}案处理。为了更准确的量刑。
        \begin{enumerate}
          \item \zd{数罪、数人存关联，共犯还有其他罪}。
          \item \al{甲法院丁的 A 罪，乙法院发现丁的 B 罪}。
                \begin{enumerate}
                  \item （非同种犯罪）一般可以并。\zd{A = B（同种犯罪）时一般应当并}。例外。If(导致过分迟延)。
                  \item 先到先得。最初受理地 \hz 主要犯罪地。
                \end{enumerate}
          \item \bs{二审并案应当发回一审重审}。
          \item \bs{网络犯罪案件。可以先并后分}。公安可以并案管辖，侦查完毕后分开移送给检察院。新增。
        \end{enumerate}
\end{enumerate}
